% Preâmbulo do livro dos tutoriais (incluído pelo pandoc em header-includes).
\usepackage{braket}      % \ket{}, \bra{}, \braket{} (notação de Dirac)
\usepackage{adjustbox}   % reduz matrizes largas para caberem na página
\usepackage{fvextra}     % quebra de linhas longas em códigos e saídas

% Códigos destacados pelo pandoc: fundo cinza-claro, fonte menor e quebra de linhas longas.
\usepackage{framed}
\definecolor{shadecolor}{HTML}{F3F1F8}
\renewenvironment{Shaded}{\begin{snugshade}}{\end{snugshade}}
\DefineVerbatimEnvironment{Highlighting}{Verbatim}{%
  commandchars=\\\{\},breaklines,breakanywhere,fontsize=\small}

% Saídas das células (texto impresso, matrizes, erros).
\definecolor{corsaida}{HTML}{8C8C8C}
\DefineVerbatimEnvironment{saida}{Verbatim}{%
  breaklines,breakanywhere,fontsize=\footnotesize,frame=leftline,framerule=1.2pt,
  rulecolor=\color{corsaida},xleftmargin=0.8em,formatcom=\color{black!75}}
\DefineVerbatimEnvironment{saidaerro}{Verbatim}{%
  breaklines,breakanywhere,fontsize=\footnotesize,frame=leftline,framerule=1.2pt,
  rulecolor=\color{red!60!black},xleftmargin=0.8em,formatcom=\color{red!60!black}}

% Cabeçalhos dos capítulos mais compactos.
\usepackage{titlesec}
\titleformat{\chapter}[display]{\normalfont\Huge\bfseries\color{purple!70!black}}{}{0pt}{}
\titlespacing*{\chapter}{0pt}{-20pt}{24pt}

% URLs longas podem quebrar em qualquer ponto.
\usepackage{xurl}

% Cabeçalho: título do capítulo em fonte menor (sem caixa alta) e número da página.
\usepackage{fancyhdr}
\pagestyle{fancy}
\fancyhf{}
\fancyhead[L]{\small\itshape\nouppercase{\leftmark}}
\fancyhead[R]{\small\thepage}
\renewcommand{\headrulewidth}{0.4pt}
\setlength{\headheight}{14pt}

% Mais folga para justificar parágrafos com nomes de funções longos (\texttt{...}).
\setlength{\emergencystretch}{3em}

% Símbolos Unicode que o pdfLaTeX não conhece (aparecem em código e em texto). Só no pdfLaTeX:
% o XeLaTeX e o LuaLaTeX já leem Unicode direto.
\ifPDFTeX
  \DeclareUnicodeCharacter{03C8}{\ensuremath{\psi}}
  \DeclareUnicodeCharacter{03A6}{\ensuremath{\Phi}}
  \DeclareUnicodeCharacter{03B1}{\ensuremath{\alpha}}
  \DeclareUnicodeCharacter{03B8}{\ensuremath{\theta}}
  \DeclareUnicodeCharacter{03BA}{\ensuremath{\kappa}}
  \DeclareUnicodeCharacter{2297}{\ensuremath{\otimes}}
  \DeclareUnicodeCharacter{221A}{\ensuremath{\surd}}
  \DeclareUnicodeCharacter{27E8}{\ensuremath{\langle}}
  \DeclareUnicodeCharacter{27E9}{\ensuremath{\rangle}}
  \DeclareUnicodeCharacter{2020}{\ensuremath{\dagger}}
\fi
